\chapter{FFT 频谱分析} \label{ch:fft}

时域波形告诉我们信号“随时间怎样变化”，却不容易直接看出信号由哪些频率组成。快速傅里叶变换（Fast Fourier Transform，FFT）把离散采样序列转换到频域，使电机振动、音频、交流电压和传感器周期干扰中的频率成分一目了然。本章先从正弦函数和离散傅里叶变换讲起，再实现基二时间抽取 FFT，最后讨论窗函数、频谱泄漏和幅值换算。完整示例代码位于 \texttt{Chapter-fft/FFT变换/} 目录。

\medskip 本章涉及到的实验内容：
\begin{enumerate}
\item 单一正弦信号与多频率叠加信号；
\item 离散傅里叶变换与基二 FFT；
\item 位逆序与蝶形运算；
\item 矩形窗、汉宁窗、海明窗和布莱克曼窗；
\item 单边幅值谱与实际频率换算。
\end{enumerate}

% ---------------------------------------------------------------------
\section{从时域观察正弦函数}

频率为 $f$、幅值为 $A$、初相位为 $\varphi$ 的正弦函数可以写成：

\begin{equation}
x(t)=A\sin(2\pi ft+\varphi)
\label{eq:sine-function}
\end{equation}

图~\ref{fig:fft-sine-functions}~依次画出了单一正弦函数、多频率叠加函数和含噪函数。单一正弦波容易从周期估计频率；多个频率叠加以后，时域波形已经很难直接读懂；加入噪声以后，周期性特征会进一步被掩盖。

\begin{figure}[htbp]
\centering
\begin{tikzpicture}[x=1.25cm,y=0.72cm]
    \draw[->,thin] (0,5.4) -- (7,5.4) node[right] {$t$};
    \draw[->,thin] (0,4.1) -- (0,6.7) node[above] {$x_1(t)$};
    \draw[structurecolor,thick] plot[domain=0:6.3,samples=120] (\x,{5.4+0.8*sin(deg(2*\x))});
    \node at (3.4,6.6) {单一正弦函数};

    \draw[->,thin] (0,2.4) -- (7,2.4) node[right] {$t$};
    \draw[->,thin] (0,1.1) -- (0,3.7) node[above] {$x_2(t)$};
    \draw[red,thick] plot[domain=0:6.3,samples=160] (\x,{2.4+0.65*sin(deg(2*\x))+0.3*sin(deg(6*\x))});
    \node at (3.4,3.6) {多频率叠加函数};

    \draw[->,thin] (0,-0.6) -- (7,-0.6) node[right] {$t$};
    \draw[->,thin] (0,-1.9) -- (0,0.7) node[above] {$x_3(t)$};
    \draw[green!50!black,thick] plot[domain=0:6.3,samples=160] (\x,{-0.6+0.65*sin(deg(2*\x))+0.3*sin(deg(6*\x))+0.14*sin(deg(31*\x))+0.1*cos(deg(43*\x))});
    \node at (3.4,0.6) {含噪函数};
\end{tikzpicture}
\caption{三类常见时域函数图} \label{fig:fft-sine-functions}
\end{figure}

% ---------------------------------------------------------------------
\section{从 DFT 到 FFT}

对长度为 $N$ 的离散序列 $x[n]$，离散傅里叶变换（Discrete Fourier Transform，DFT）定义为：

\begin{equation}
X[k]=\sum_{n=0}^{N-1}x[n]\exp\left(-\mathrm{j}\frac{2\pi kn}{N}\right),\quad 0\leq k<N
\label{eq:dft}
\end{equation}

直接按照定义计算，每个频点都要遍历全部采样点，总计算量为 $O(N^2)$。FFT 不是另一种变换，而是利用旋转因子的周期性和对称性，加速 DFT 的计算方法。基二 FFT 把偶数下标和奇数下标数据不断拆分，使计算量降为 $O(N\log_2N)$。

\begin{anybox}{采样点数要求}
本章实现的是基二时间抽取 FFT，因此采样点数必须是二的整数次幂，例如 64、128、256 和 1024。不是二的整数次幂时，接口会返回参数错误。
\end{anybox}

% ---------------------------------------------------------------------
\section{位逆序与蝶形运算}

\subsection{位逆序重排}

以八点 FFT 为例，下标 $3$ 的二进制是 $011$，反转三位以后得到 $110$，即十进制 $6$。位逆序把原始顺序重排为 $0$、$4$、$2$、$6$、$1$、$5$、$3$ 和 $7$，随后每一级蝶形都能在数组中原地完成。

\subsection{蝶形运算}

一个基二蝶形接收复数 $a$ 与 $b$，乘以旋转因子 $W_N^k$ 后生成两个输出：

\begin{align}
X_0 &= a+W_N^kb \\
X_1 &= a-W_N^kb \\
W_N^k &= \exp\left(-\mathrm{j}\frac{2\pi k}{N}\right)
\label{eq:fft-butterfly}
\end{align}

图~\ref{fig:fft-butterfly}~中的交叉结构像蝴蝶翅膀，因此得名蝶形运算。

\begin{figure}[htbp]
\centering
\begin{tikzpicture}[
    point/.style={circle,fill=structurecolor,inner sep=2pt},
    arrow/.style={-stealth,thick}
]
    \node (a) at (-4,1.2) {$a$};
    \node (b) at (-4,-1.2) {$b$};
    \node[point] (p1) at (-1.0,1.2) {};
    \node[point] (p2) at (-1.0,-1.2) {};
    \node (x0) at (3.5,1.2) {$a+W_N^kb$};
    \node (x1) at (3.5,-1.2) {$a-W_N^kb$};

    \draw[arrow] (a) -- (p1) -- (x0);
    \draw[arrow] (b) -- node[below] {$W_N^k$} (p2) -- (x1);
    \draw[arrow] (p1) -- node[pos=0.55,above] {$+$} (x1);
    \draw[arrow] (p2) -- node[pos=0.55,below] {$+$} (x0);
    \node at (1.0,-0.15) {$-$};
\end{tikzpicture}
\caption{基二 FFT 蝶形运算图} \label{fig:fft-butterfly}
\end{figure}

八点 FFT 需要三层蝶形。一般地，$N$ 点基二 FFT 需要 $\log_2N$ 层，每一层处理 $N$ 个数据。图~\ref{fig:fft-flow}~概括了程序中的数据流。

\begin{figure}[htbp]
\centering
\begin{tikzpicture}[
    block/.style={draw=structurecolor, thick, rounded corners=2pt, minimum width=2.5cm, minimum height=1.0cm, align=center, fill=structurecolor!8},
    arrow/.style={-stealth,thick}
]
    \node[block] (sample) at (-5.0,0) {定时采样};
    \node[block] (window) at (-2.0,0) {乘窗函数};
    \node[block] (reverse) at (1.0,0) {位逆序重排};
    \node[block] (butterfly) at (4.0,0) {逐层蝶形};
    \node[block] (magnitude) at (4.0,-2.0) {单边幅值谱};

    \draw[arrow] (sample) -- (window);
    \draw[arrow] (window) -- (reverse);
    \draw[arrow] (reverse) -- (butterfly);
    \draw[arrow] (butterfly) -- (magnitude);
\end{tikzpicture}
\caption{FFT 频谱分析流程图} \label{fig:fft-flow}
\end{figure}

% ---------------------------------------------------------------------
\section{频点、分辨率与单边频谱}

采样频率为 $f_s$，采样点数为 $N$ 时，频率分辨率为：

\begin{equation}
\Delta f=\frac{f_s}{N}
\label{eq:fft-resolution}
\end{equation}

第 $k$ 个频点对应的实际频率为：

\begin{equation}
f_k=k\Delta f=k\frac{f_s}{N}
\label{eq:fft-bin}
\end{equation}

例如采样频率为 $128\,\mathrm{Hz}$，数据长度为 $128$，频率分辨率就是 $1\,\mathrm{Hz}$。第八个频点对应 $8\,\mathrm{Hz}$，第二十个频点对应 $20\,\mathrm{Hz}$。

实数时域信号的频谱关于零频率共轭对称，工程中通常只显示直流到奈奎斯特频率之间的单边频谱。除直流分量以外，正频率幅值需要乘二：

\begin{equation}
A[k]=\begin{cases}
\dfrac{|X[0]|}{N}, & k=0 \\[6pt]
\dfrac{2|X[k]|}{N}, & 1\leq k<\dfrac{N}{2}
\end{cases}
\end{equation}

% ---------------------------------------------------------------------
\section{窗函数与频谱泄漏}

FFT 把有限长度数据当作一个周期不断重复。如果截取的数据刚好包含整数个信号周期，首尾可以平滑连接；如果不是整数周期，首尾会出现跳变，能量就会扩散到相邻频点，这种现象称为频谱泄漏。乘窗函数可以压低数据两端，减小这个跳变。

\subsection{四种常用窗函数}

矩形窗不改变原始数据：

\begin{equation}
w_{R}[n]=1
\end{equation}

汉宁窗、海明窗和布莱克曼窗分别为：

\begin{align}
w_{Hn}[n] &= 0.5-0.5\cos\left(\frac{2\pi n}{N-1}\right) \\
w_{Hm}[n] &= 0.54-0.46\cos\left(\frac{2\pi n}{N-1}\right) \\
w_{B}[n] &= 0.42-0.5\cos\left(\frac{2\pi n}{N-1}\right)
            +0.08\cos\left(\frac{4\pi n}{N-1}\right)
\label{eq:fft-windows}
\end{align}

图~\ref{fig:fft-window-functions}~显示了四种窗函数的形状。横轴使用归一化采样位置，因此不依赖具体点数。

\begin{figure}[htbp]
\centering
\begin{tikzpicture}[x=10cm,y=4.5cm]
    \draw[->,thick] (0,0) -- (1.08,0) node[right] {归一化位置};
    \draw[->,thick] (0,0) -- (0,1.18) node[above] {权重};
    \draw[gray,thick] (0,1) -- (1,1);
    \draw[structurecolor,thick] plot[domain=0:1,samples=100] (\x,{0.5-0.5*cos(deg(2*pi*\x))});
    \draw[red,thick] plot[domain=0:1,samples=100] (\x,{0.54-0.46*cos(deg(2*pi*\x))});
    \draw[green!50!black,thick] plot[domain=0:1,samples=100] (\x,{0.42-0.5*cos(deg(2*pi*\x))+0.08*cos(deg(4*pi*\x))});
    \node[gray] at (0.83,1.08) {矩形窗};
    \node[structurecolor] at (0.78,0.72) {汉宁窗};
    \node[red] at (0.78,0.58) {海明窗};
    \node[green!50!black] at (0.78,0.43) {布莱克曼窗};
\end{tikzpicture}
\caption{四种常用窗函数图} \label{fig:fft-window-functions}
\end{figure}

矩形窗的主瓣最窄，频率分辨能力强，但旁瓣高；布莱克曼窗的旁瓣低，抑制泄漏能力强，但主瓣更宽；汉宁窗和海明窗位于两者之间。没有一种窗函数在所有指标上都最好，选择取决于更关心“分开两个接近频率”，还是更关心“看见强信号旁边的弱信号”。

\subsection{加窗前后的频谱}

图~\ref{fig:fft-leakage-comparison}~以频率没有落在整数频点上的正弦信号为例。矩形窗下能量分散到较宽范围，汉宁窗降低了远端旁瓣，但主峰附近变宽。

\begin{figure}[htbp]
\centering
\begin{tikzpicture}[x=0.45cm,y=3.6cm]
    \draw[->,thick] (0,0) -- (13.2,0) node[right] {频点};
    \draw[->,thick] (0,0) -- (0,1.25) node[above] {幅值};
    \draw[gray,thick] plot[smooth] coordinates {(0,0.03) (1,0.04) (2,0.05) (3,0.07) (4,0.10) (5,0.16) (6,0.30) (7,0.86) (8,0.66) (9,0.25) (10,0.14) (11,0.09) (12,0.06)};
    \draw[structurecolor,thick] plot[smooth] coordinates {(0,0.00) (1,0.00) (2,0.00) (3,0.01) (4,0.02) (5,0.08) (6,0.42) (7,0.92) (8,0.78) (9,0.28) (10,0.04) (11,0.01) (12,0.00)};
    \node[gray] at (10.2,1.05) {矩形窗};
    \node[structurecolor] at (10.2,0.88) {汉宁窗};
\end{tikzpicture}
\caption{加窗前后的频谱泄漏对比} \label{fig:fft-leakage-comparison}
\end{figure}

窗函数会改变幅值。需要精确测量幅值时，应根据所选窗函数的相干增益做补偿。本章示例重点演示频率成分识别，因此只完成标准单边幅值换算，没有额外加入窗增益补偿。

% ---------------------------------------------------------------------
\section{时域与频域对照}

设测试信号由 $8\,\mathrm{Hz}$ 和 $20\,\mathrm{Hz}$ 两个正弦分量组成：

\begin{equation}
x(t)=\sin(2\pi\cdot8t)+0.5\sin(2\pi\cdot20t)
\end{equation}

图~\ref{fig:fft-time-frequency}~左侧时域波形比较复杂，右侧频域却只在两个频率附近出现明显峰值。幅值较小的二十赫兹分量，对应频谱中的峰值也较低。

\begin{figure}[htbp]
\centering
\begin{tikzpicture}[x=0.72cm,y=1.4cm]
    \draw[->,thick] (0,0) -- (7.2,0) node[right] {时间};
    \draw[->,thick] (0,-1.5) -- (0,1.7) node[above] {幅值};
    \draw[structurecolor,thick] plot[domain=0:6.3,samples=180] (\x,{sin(deg(2*\x))+0.5*sin(deg(5*\x))});
    \node at (3.3,-1.75) {时域波形};

    \begin{scope}[xshift=9cm]
        \draw[->,thick] (0,0) -- (7.2,0) node[right] {频率};
        \draw[->,thick] (0,0) -- (0,1.7) node[above] {幅值};
        \draw[structurecolor,very thick] (2,0) -- (2,1.25);
        \draw[red,very thick] (5,0) -- (5,0.63);
        \node[below] at (2,0) {$8\,\mathrm{Hz}$};
        \node[below] at (5,0) {$20\,\mathrm{Hz}$};
        \node at (3.4,-1.25) {单边幅值谱};
    \end{scope}
\end{tikzpicture}
\caption{双频信号的时域与频域对照图} \label{fig:fft-time-frequency}
\end{figure}

% ---------------------------------------------------------------------
\section{程序实现}

示例代码提供五个公共函数：

\begin{itemize}
\item \texttt{FFT\_IsPowerOfTwo}：检查数据长度；
\item \texttt{FFT\_ApplyWindow}：施加四种窗函数之一；
\item \texttt{FFT\_Transform}：执行原地基二时间抽取 FFT；
\item \texttt{FFT\_CalcMagnitude}：计算单边幅值谱；
\item \texttt{FFT\_GetBinFrequency}：把频点下标换算为实际频率。
\end{itemize}

典型调用流程如下：

\begin{lstlisting}[style=c]
/* ADC 采样值放入实部，虚部清零 */
FFT_ApplyWindow(realPart, sampleCount, FFT_WINDOW_HANN);
FFT_Transform(realPart, imaginaryPart, sampleCount);
FFT_CalcMagnitude(realPart, imaginaryPart, magnitude, sampleCount);

/* 将峰值下标换算成实际频率 */
peakFrequency = FFT_GetBinFrequency(peakIndex,
                                    sampleFrequency,
                                    sampleCount);
\end{lstlisting}

FFT 包含较多乘法、三角函数和浮点数组。STM32 移植时应关注三点：采样必须由定时器触发，避免主循环抖动；数组较大时应放在静态存储区，避免压垮任务栈；芯片已经提供硬件浮点单元或 CMSIS-DSP 时，可以用经过优化的库函数替换教学实现。

% ---------------------------------------------------------------------
\section{PC 验证结果}

本章示例使用 C99 编译，并打开 \texttt{-Wall}、\texttt{-Wextra}、\texttt{-Wpedantic} 和 \texttt{-Werror} 严格警告选项。链接时加入数学库，\texttt{fft\_transform.c} 与 \texttt{main.c} 编译过程没有产生警告，程序正常结束。

测试信号由幅值为 $1$ 的 $8\,\mathrm{Hz}$ 正弦波和幅值为 $0.5$ 的 $20\,\mathrm{Hz}$ 正弦波组成。采样频率为 $128\,\mathrm{Hz}$，FFT 点数为 128，因此频率分辨率为 $1\,\mathrm{Hz}$。程序输出的两个主要峰值如表~\ref{tab:fft-pc-result}~所示。

\begin{table}[htbp]
\centering
\caption{FFT 示例的 PC 频谱结果} \label{tab:fft-pc-result}
\begin{tabular}{llll}
\toprule
频点下标 & 实际频率 & 理论幅值 & 计算幅值 \\
\midrule
8 & $8.00\,\mathrm{Hz}$ & 1.0000 & 1.0000 \\
20 & $20.00\,\mathrm{Hz}$ & 0.5000 & 0.5000 \\
\bottomrule
\end{tabular}
\end{table}

结果表明频点换算与单边幅值归一化符合测试信号。独立断言还验证了二次幂长度判断、汉宁窗端点、海明窗端点、布莱克曼窗端点和八点单频 FFT，断言均正常通过。

% ---------------------------------------------------------------------
\section{完整示例代码}

完整示例位于 \texttt{Chapter-fft/FFT变换/}：

\begin{itemize}
\item \texttt{fft\_transform.h}：窗函数枚举与 FFT 公共函数；
\item \texttt{fft\_transform.c}：位逆序、蝶形、窗函数和幅值计算；
\item \texttt{main.c}：双频测试信号生成与频谱输出。
\end{itemize}

% ---------------------------------------------------------------------
\section{本章小结}

FFT 的输出不是抽象的数学结果，而是一组有明确物理含义的频点。采样频率决定可观察的最高频率，采样点数决定频率分辨率，窗函数决定主瓣宽度和旁瓣高度。只有同时理解采样、加窗、变换和幅值换算，频谱图上的峰值才能正确对应到真实系统中的振动、噪声与周期信号。

% ---------------------------------------------------------------------
\begin{problemset}
\item 采样频率为 $1024\,\mathrm{Hz}$、FFT 点数为 256 时，频率分辨率和不发生混叠时可观察的最高频率分别是多少？下标为 37 的频点对应多少赫兹？
\item 对由 $60\,\mathrm{Hz}$ 和 $180\,\mathrm{Hz}$ 两个正弦分量组成的信号进行 256 点 FFT，采样频率为 $1024\,\mathrm{Hz}$。两个分量理论上分别落在哪个频点？它们是否恰好落在整数频点上？
\item 什么是频谱泄漏？比较矩形窗与汉宁窗时，为什么不能简单地认为“旁瓣更低的窗函数在任何情况下都更好”？
\item 基二时间抽取 FFT 为什么要求输入长度是二的整数次幂？位逆序重排和蝶形运算在算法中分别承担什么作用？
\item 在 STM32 上使用定时器触发 ADC、DMA 双缓冲采集 512 点数据。请说明 DMA 中断和频谱处理任务各自应该完成什么工作，并给出判断系统能否实时运行的时间约束。
\end{problemset}
